\documentclass[10pt,a4paper]{amsart}
\usepackage[margin=19mm]{geometry}
\usepackage{amsmath,amssymb,mathtools}
\usepackage[scr=rsfs,cal=euler]{mathalfa}
\usepackage{xfrac}
\usepackage[unicode,hidelinks,bookmarks=false]{hyperref}
\newcommand{\ie}{i.e.\ }
\newcommand{\cf}{cf.\ }
\newcommand{\resp}{respectively\ }
\newcommand{\Subsection}{\subsection}
\providecommand{\nobreakdash}{\nobreak\hbox{-}}
\setlength{\emergencystretch}{2em}
\multlinegap0pt
\usepackage{amssymb}
\usepackage{mathtools}
\usepackage[shortlabels]{enumitem}
\usepackage{float}
\usepackage{xcolor}
\usepackage{tikz}
\newtheorem{definition}{Definition}
\newtheorem{lemma}{Lemma}
\newtheorem{corollary}{Corollary}
\newtheorem{theorem}{Theorem}
\title{The Erd\H{o}s--Hajnal Property for the six-vertex Graph\\ with Edge Set $\{ab,bc,cd,de,af,bf,df\}$}
\author{Viet-Hoang Tran, Tan M. Nguyen}
\date{Source: arXiv:2608.28551v1 (CC BY 4.0)}
\begin{document}
\maketitle

\noindent\emph{Source and licence.} The Erd\H{o}s--Hajnal Property for the six-vertex Graph\\ with Edge Set $\{ab,bc,cd,de,af,bf,df\}$. Version arXiv:2608.28551v1 (CC BY 4.0), distributed by arXiv under the Creative Commons Attribution 4.0 International licence, http://creativecommons.org/licenses/by/4.0/, as recorded in the arXiv licence metadata for this exact version and shown on the abstract page https://arxiv.org/abs/2608.28551v1. Full licence notice: \texttt{licenses/arxiv-2608.28551-CC-BY-4.0.txt}.\par\smallskip

\noindent\textit{Editorial scope.} Counted unit: the article's theorem thm:F-comb-conclusion (source0.tex lines 3387--3391). The article's complete original proof of it is delivered with it (source0.tex lines 3393--3512, one \texttt{proof} environment of 120 lines). No mathematics is written, repaired or re-derived: the statement and the proof are the article's own lines, in the article's own notation, and no retained line is substituted or re-flowed.  Retained uncounted as the source-local context the proof needs: lines 468--487; lines 1806--1826; lines 3188--3198; lines 3252--3274; lines 3320--3342; lines 3353--3359.  Nothing was omitted from any retained span, so the package carries no omission record.  No retained line contains a citation, so the wrapper carries no bibliography and no bibliography entry.  No retained line contains a figure environment, an image-inclusion command or drawing code, so no image is delivered and the package declares no asset dependency.  Editorial formatting only, in addition: an AMS \texttt{amsart} preamble that restates the article's own declarations and macros used by the retained lines, namely the environments \texttt{definition}, \texttt{lemma}, \texttt{corollary}, \texttt{theorem} on separate counters, together with the standard packages those lines need.  The retained environments print in this document as Definition 1, Lemma 1, Corollary 1, Theorem 1 in order of appearance, whereas the article prints the same items with the numbering of its own full document.  The complete native source of the article as distributed in the arXiv e-print for this exact version is bundled unchanged as \texttt{sources/208795/source0.tex}.
For a positive integer $\ell$ and a real number $w>0$, an
\emph{$(\ell,w)$-comb} is a sequence
\[
                     ((h_i,B_i):i\in[\ell])
\]
such that $(B_1,\ldots,B_\ell)$ is an $(\ell,w)$-blockade, the
vertices $h_1,\ldots,h_\ell$ are distinct and disjoint from all the
blocks, and, for distinct $i,j$, the vertex $h_i$ is complete to
$B_i$ and anticomplete to $B_j$.  We call the $h_i$ the
\emph{handles} and the $B_i$ the \emph{teeth}.  Property $(*)$ below
also assumes a vertex
\begin{equation}
 \begin{split}
 r&\notin\{h_1,\ldots,h_\ell\}\cup\bigcup_iB_i,\\
 r&\text{ is complete to every tooth and anticomplete to every handle.}
 \end{split}
 \label{eq:root-outside}
\end{equation}
We will always say explicitly that the comb is equipped with such a
vertex $r$.

\begin{definition}[Quantitative comb condition]
\label{def:comb-hypothesis}
A finite graph family $\mathcal F$ satisfies the
\emph{quantitative comb condition} if there are constants
\[
                       a,b,g,p,q>0
\]
such that every $(\ell,w)$-comb in every
$\overline{\mathcal F}$-free graph $G$, with $\ell,w\geq4$, equipped with
a vertex $r$ satisfying \eqref{eq:root-outside}, has one of the
following outcomes:
\begin{enumerate}
\item[(H)] $\operatorname{hom}(G)\geq w^a$;
\item[(U)] a complete or anticomplete $(k,w/k^b)$-blockade for some
           $k\geq\ell^g$;
\item[(P)] a pure $(k,w/\ell^q)$-blockade for some
           $\ell^p\leq k\leq\ell$.
\end{enumerate}
After replacing $p$ by $\min\{p,1\}$, we may and shall assume that
$0<p\leq1$.
\end{definition}

\begin{lemma}[Pivot--true-twin extraction]
\label{lem:pivot-twin-extraction}
Let $H$ be the graph induced by $\ell$ handles, where $\ell\geq16$.
At least one of the following holds.
\begin{enumerate}
 \item $H$ has a clique of order at least $\ell^{1/4}$;
 \item there is a set $I$ of at least $\ell^{1/4}/2$ handle indices
       such that, whenever $h_ih_j$ is an edge with $i,j\in I$, the
       whole interface $B_i$--$B_j$ is anticomplete.
\end{enumerate}
\end{lemma}

\begin{lemma}[Components and anticomponents]
\label{lem:pivot-carriers}
Let teeth $B_i$, $i\in I$, have order at least $w$, and let $L\geq2$
be an integer.
One of the following holds.
\begin{enumerate}
 \item $G$ has a complete or anticomplete blockade of length $L$ and
       width at least $w/L^3$;
 \item for every $i\in I$ there are
       $X_i\subseteq C_i\subseteq B_i$ such that $G[C_i]$ is connected,
       $G[X_i]$ is anticonnected, and
       \[
                         |X_i|\geq w/L^2.
       \]
\end{enumerate}
If, in addition, every handle edge on $I$ has an anticomplete tooth
interface, then in the second outcome every $X_i$--$X_j$ matrix is additively separable over $\mathbb F_2$:
\begin{equation}
 M_{ij}(x,y)=\lambda_i^j(x)\oplus\lambda_j^i(y)
                                      \oplus\gamma_{ij}.
 \label{eq:pivot-carrier-affine}
\end{equation}
\end{lemma}

\begin{corollary}[Extraction from additively separable interfaces]
\label{cor:pivot-affine-extraction}
There are constants $\rho,\sigma>0$ and an integer $m_0$ with the
following property.  In a $T$-free graph containing a comb and a vertex
complete to all teeth and anticomplete to all handles, let
$X_1,\ldots,X_m$ be anticonnected subsets of distinct
teeth, where $m\geq m_0$ and $|X_i|\geq W$ for every $i$.  Suppose
every pairwise interface has the form
\[
 M_{ij}(x,y)=\lambda_i^j(x)\oplus\lambda_j^i(y)
                         \oplus\gamma_{ij},
\]
and suppose $h_ih_j$ is a nonedge whenever the $X_i$--$X_j$ interface
is mixed.  Then one of the following holds:
\begin{enumerate}
 \item a pure blockade of length at least $m^\rho$ and width at least
       $W$;
 \item a complete or anticomplete blockade of length at least
       $m^\sigma$ and width at least $W/m$.
\end{enumerate}
In either outcome, the blocks lie in pairwise distinct sets $X_i$; in
particular, the actual length of the blockade is at most $m$.
\end{corollary}

\begin{lemma}[Clique-handle extraction]
\label{lem:pivot-clique-extraction}
For $m\geq144$ clique handles with teeth of order at least $w$, there
is a complete or anticomplete blockade of length at least $\sqrt m/2$
and width at least $w/(2m)$.  For every $m\geq3$ there is a
pure pair of width $w/(4m)$.
\end{lemma}

\begin{theorem}
\label{thm:F-comb-conclusion}
The singleton family $\{F\}$ satisfies the quantitative comb condition of
Definition~\ref{def:comb-hypothesis}.
\end{theorem}

\begin{proof}
Let $\rho,\sigma>0$ and $m_0$ be the constants in
Corollary~\ref{cor:pivot-affine-extraction}.  Decreasing $\rho$ if
necessary, assume that $0<\rho\leq1$.  Set
\[
 q_0=\frac{1}{8},\qquad p_0=\frac{\rho}{8},\qquad
 c_*=\min\left\{\frac{1}{16},\frac{\sigma}{8}\right\}.
\]
Choose an integer $\ell_0$ so large that, whenever $\ell\geq\ell_0$,
all of the following hold:
\begin{enumerate}
 \item $L=\lfloor\ell^{q_0}\rfloor\geq2$ and
       $L\geq\ell^{1/16}$;
 \item $\ell^{1/4}/2\geq m_0$;
 \item every clique of at least $\ell^{1/4}$ handles has order at
       least $144$;
 \item
 \[
  \frac{1}{2}\ell^{1/8}\geq\ell^{c_*},\qquad
  2^{-\rho}\ell^{\rho/4}\geq\ell^{p_0},\qquad
  2^{-\sigma}\ell^{\sigma/4}\geq\ell^{c_*};
  \]
\end{enumerate}

After fixing $\ell_0$, put
\[
 c=\min\left\{c_*,\frac{\log 2}{\log\ell_0}\right\},
 \qquad p=p_0,\qquad q=\frac{1}{4},
\]
and choose
\begin{equation}
 b\geq\max\left\{3,\frac{2}{c},
          \log_2\bigl(\max\{8,4\ell_0\}\bigr)\right\}.
 \label{eq:comb-b-choice}
\end{equation}
We verify Definition~\ref{def:comb-hypothesis} with
$a=1$, $g=c$, and these fixed $b,p,q$.  Outcome~(H) will not be needed.

Consider first $\ell\geq\ell_0$.  Apply
Lemma~\ref{lem:pivot-twin-extraction}.  Suppose it gives a clique of
$m\geq\ell^{1/4}$ handles.  By
Lemma~\ref{lem:pivot-clique-extraction}, there is a complete or
anticomplete blockade of length at least
$\sqrt m/2\geq\ell^{1/8}/2$ and width at least
\[
                         \frac{w}{2m}\geq \frac{w}{2\ell}.
\]
Keep exactly $k=\lceil\ell^c\rceil$ of its blocks.  By the choice of
$\ell_0$, its integer length is at least
$\lceil\ell^{c_*}\rceil\geq k$, so this is possible; also
$k\geq\ell^c$.  Since
$b c\geq2$, we have $k^b\geq\ell^{bc}\geq\ell^2\geq2\ell$, and
therefore the retained width is at least $w/k^b$.  This is outcome (U).

We may thus assume that Lemma~\ref{lem:pivot-twin-extraction} gives a
set $I$ satisfying
\[
                         |I|\geq\ell^{1/4}/2,
\]
and such that every handle edge on $I$ has an anticomplete tooth
interface.  Apply Lemma~\ref{lem:pivot-carriers} with
$L=\lfloor\ell^{1/8}\rfloor$.  If its first conclusion holds, it
gives a complete or anticomplete blockade with $k=L\geq\ell^c$ blocks
and width at least $w/L^3\geq w/k^b$, because $b\geq3$.  This is again
outcome (U).

In the other conclusion, there are
$X_i\subseteq C_i\subseteq B_i$, $i\in I$, such that $G[C_i]$ is
connected, $G[X_i]$ is anticonnected, and
\[
                         |X_i|\geq \frac{w}{L^2}
                                  \geq \frac{w}{\ell^{1/4}}.
\]
Every interface $X_i$--$X_j$ is additively separable over
$\mathbb F_2$.  On a handle edge the interface is constant zero; hence
every mixed interface lies over a handle nonedge.  We may therefore
apply Corollary~\ref{cor:pivot-affine-extraction} with
$m=|I|$ and $W=w/\ell^{1/4}$.

In its pure conclusion, the length $k$ satisfies
\[
 k\geq |I|^\rho
       \geq2^{-\rho}\ell^{\rho/4}
       \geq\ell^p,
 \qquad k\leq |I|\leq\ell,
\]
and the width is at least $w/\ell^{1/4}=w/\ell^q$.  This is outcome
(P).

In its complete-or-anticomplete conclusion, the length is at least
\[
 |I|^\sigma\geq2^{-\sigma}\ell^{\sigma/4}
                   \geq\ell^{c_*}\geq\ell^c,
\]
and the width is at least
\[
 \frac{W}{|I|}\geq \frac{w}{\ell^{5/4}}.
\]
Keep $k=\lceil\ell^c\rceil$ blocks.  These blocks are available because
the original integer length is at least
$\lceil\ell^{c_*}\rceil\geq k$.  Also $k^b\geq\ell^2$, so
$w/\ell^{5/4}\geq w/k^b$.  This is outcome (U).  This proves the
required conclusion for $\ell\geq\ell_0$.

It remains to consider $4\leq\ell<\ell_0$.  If two handles are
nonadjacent, apply Lemma~\ref{lem:pivot-carriers} with $L=2$ to their
two teeth.  Its first conclusion gives a complete or anticomplete pair
of width at least $w/8$.  In its second conclusion the two retained
sets have order at least $w/4$, their interface is additively separable,
and a largest fibre of each of the two binary endpoint functions gives
a complete or anticomplete pair of width at least $w/8$.  If every pair
of handles is adjacent, Lemma~\ref{lem:pivot-clique-extraction} gives a
complete or anticomplete pair of width at least $w/(4\ell_0)$.

In every bounded case take $k=2$.  The definition of $c$ gives
$\ell^c\leq2$, while the choice of $b$ in
\eqref{eq:comb-b-choice} gives
$2^b\geq\max\{8,4\ell_0\}$.  Thus the pair has width at least
$w/2^b$, and outcome (U) holds.  The theorem follows.
\end{proof}

\end{document}
